% Build through paper/build.py; frozen dopamine values are bound in result-bindings.json.
\documentclass[11pt]{article}
\usepackage[margin=1in]{geometry}
\usepackage{graphicx,booktabs,microtype,xcolor,hyperref}
\usepackage{caption,xfp,siunitx}
\captionsetup{justification=raggedright,singlelinecheck=false}
\definecolor{pending}{HTML}{9B3D14}
\newif\iffinal
\input{final-mode.tex}
\newcommand{\paperstatus}{PREPRINT DRAFT --- HUMAN AUTHOR APPROVAL PENDING}
\newcommand{\result}[1]{\textcolor{pending}{[UNFILLED RESULT]}}
\newcommand{\fixedthree}[1]{\num[round-mode=places,round-precision=3,round-pad=true]{#1}}\def\roundedbounds[#1, #2]\relax{[\(\fixedthree{#1}\), \(\fixedthree{#2}\)]}
\newcommand{\roundedci}[1]{\expandafter\roundedbounds#1\relax}\def\roundedmixedbounds[#1, #2]\relax{[\(\fixedthree{#1}\), \(+\fixedthree{#2}\)]}\newcommand{\roundedmixedci}[1]{\expandafter\roundedmixedbounds#1\relax}
\title{A whole-nervous-system model of the male fly as a pharmacology bench}
\author{Hasan ``Rolf'' Yildirim\\Lasell University}
\date{}\hypersetup{pdftitle={A whole-nervous-system model of the male fly as a pharmacology bench},pdfauthor={Hasan "Rolf" Yildirim}}
\begin{document}
\maketitle
\iffinal\else
\begin{center}\small\bfseries\color{pending}\paperstatus\end{center}
\fi
\begin{abstract}
We asked whether a simulated male fly nervous system can help with pharmacology questions without confusing a model with an animal. In one fixed connectome-based model, we perturbed dopamine clearance, disconnected transmitter classes, stimulated a courtship circuit and varied inhibitory conductance. The initial dopamine history contrast was nominally significant before correction but failed its multiple-test correction, and an independently planned follow-up found no detected difference. Removing major transmitters produced large changes in the expected directions, and courtship stimulation reached song-related neurons; neither observation validates behaviour. The clearest graded pharmacology result is a late rise in firing as the model's GABA brake is blocked, partly compensated by another inhibitory channel. The virtual body only plays back recorded spikes, so it shows no behavioural outcome.
\end{abstract}

\section*{Introduction: the model and the question}
A connectome is a map of cells and connections. We simulate the MaleCNS male central nervous system \cite{berg} with point neurons that emit spikes when their model voltage crosses a threshold, following earlier connectome-constrained simulations \cite{shiu}. The retained graph has 162,517 neurons and 25,120,209 directed connections. The anatomy picture (Fig.~\ref{fig:anatomy}) shows where the cells are and contains no physiological measurement. We inject background drive into sensory cells; in the history experiments we also inject patterned drive into tubercle-to-bulb (TuBu) cells downstream of the eye. The fly does not see. The neural model has no learning, no visual transduction into these history inputs and no validated movement.

A whole-CNS model is interesting for pharmacology because a local change need not have a local outcome. Changing a synaptic brake can alter activity in cells far from those receiving the input, and a change that barely shifts one readout may be prominent in another. A connectome-based simulation lets us keep the wiring fixed, change an assumed transporter or connection strength, and ask which readout moved. Shiu and colleagues showed that a connectome-constrained fly model can be used to study sensorimotor processing and test some predictions experimentally \cite{shiu}. That precedent motivates our approach but does not validate our male model's physiology. Our model is built on the male CNS anatomy from Berg and colleagues \cite{berg}. Anatomy shows possible routes for an effect. It gives no measured dose response, and it does not guarantee that any simulated spike predicts a living fly's response.

The original motivation was an ADHD-like question. The \textit{fumin} fly mutant has impaired dopamine-transporter cleanup and is a model of hyperactivity, not an ``ADHD fly'' \cite{kume,shin}. Human ADHD, attention and locomotion are not measured here. The model's dopamine receptor densities are assumed and largely uniform. We therefore ask narrower questions: does transporter removal change a specified neural response to earlier input, and do transmitter or inhibitory-conductance changes produce interpretable model activity?

\noindent\fbox{\begin{minipage}{\dimexpr\linewidth-2\fboxsep-2\fboxrule\relax}\small
\textbf{Key terms.} \emph{Connectome}: map of neurons and connections. \emph{Spiking model}: simplified equations for electrical events. \emph{Transmitter knockout}: removal, in the simulation, of the outgoing connections assigned to one chemical class; the neurons themselves are kept. \emph{GABA}: a model inhibitory connection class; \emph{GluCl}: a separate glutamate-gated inhibitory class. \emph{Blocked fraction}: the assumed fraction of GABA conductance removed. \emph{History contrast}: a difference in neural response to the same final paired input after different earlier inputs. \emph{Seed}: one random repeat on the same anatomical network, not another animal. \emph{Bootstrap interval}: uncertainty estimated by resampling whole repeats together. An interval including zero means \emph{no detected difference}, not equality.
\end{minipage}}

\section*{Methods in brief}
Supplementary tables S1 to S5 give the model summary, population breakdown, neuron and synapse parameters, update equations and identifiers of the named cell types.

In the dopamine study, artificial TuBu input has two independently calibrated patterns, A and B, at $-50^\circ$ and $+50^\circ$. We measure the 245 ellipsoid-body ring neurons receiving direct TuBu input. We compare their responses to the same final paired pattern after subtracting responses with no final input. That difference is the history contrast $H$. One template unit equals the calibrated A-minus-B response. We compare $H$ between a normal transporter and its removal (modelled \textit{fumin}), each with and without a chemistry-fixed reduction in dopamine release set by a model multiplier. Ten repeated runs formed the initial comparison; forty new runs formed the follow-up. Each trial settled for 2 s, then used a 5 s prefix, 1 s gap and 14 s test; the first 2 s of the test were the primary window. We specified paired random repeats, an independently fitted projection axis, whole-repeat bootstrap intervals \cite{efron} and sign-randomisation tests \cite{ernst} before seeing outcomes. Two initial primary contrasts share a Holm multiple-test correction \cite{holm}; a separate follow-up froze only the genotype contrast and its decision rule in advance. Before the analysed pilot, partial runs were voided after an input-plumbing repair. The retained calibration failed its original random-number audit; only part of it was replayed. The initial affine-control calculation, meant to separate history change from simple per-cell scaling, was numerically unstable, which rules it out as mechanism evidence.

For the circuit tour, the same network receives injected background drive while the outgoing connections of one transmitter class at a time are switched off. A separate series injects low-to-high input directly into P1 courtship neurons or pIP10 descending neurons. Ten repeated runs per condition estimate firing-rate changes relative to controls, after 2 s of settling and over 10 s of measurement. Courtship runs instead used 5 s of stimulation and 5 s after stimulation; input was 10, 30 or 60 Hz at P1, and 30 Hz at pIP10 or random control cells. The tour uses only disconnections and imposed electrical stimuli.

For the GABA-response experiment, a uniform GABA-class conductance multiplier implements Rdl-like block. We call the fraction of GABA conductance removed $\theta$. Other conditions increase GABA conductance or strengthen GluCl conductance under partial GABA block; those axes are in multiples of baseline conductance. The dose design and Hill readouts were frozen before the runs; the half-point relative to the observed full-block endpoint was defined at analysis. Ten paired runs per condition settled for 2 s and were measured for 10 s. Matched whole-run bootstrap intervals describe differences in firing outside the directly driven sensory cells. These runs use dark rest with background drive injected into sensory cells only, and no TuBu stimulus; the fly does not see. No ingestion, distribution, native receptor sensitivity or off-target channel action is modelled.

The original spiking calculations used Brian2 \cite{stimberg}. A separate CUDA engine uses CuPy and NVRTC for runtime compilation and Philox for counter-based random input \cite{okuta,nvrtc,salmon}. It keeps the connection map on a GPU and advances voltage, spikes and delayed connections in a persistent kernel, updating dopamine at its model cadence. With identical saved background events and clamped dopamine, all 58,655,422 emitted spikes across four 12-second intact and GABA-disconnected trajectories matched Brian2 at the 0.1 ms output grid, with one network per run. That match tests the computational implementation alone. On an NVIDIA L4, one 12-second run took 8.420 s of stepping at the 10 ms interface, and an independently timed short run took 0.637029 wall seconds per simulated second at that interface and 1.268480 at the 1 ms interface. These are single observations of GPU stepping, not speedup estimates, and they exclude network construction, startup and a closed-loop body. We also re-ran the full GABA ladder with the same design but independent random streams, pairing each engine's effects against its own controls. Unlike identical-input spike matching, overlapping effect intervals here describe agreement in the measured curve, not spike-train or statistical equivalence \cite{alevi,knight}. Full implementation and timing details are in the separate GPU methods write-up.

\section*{Dopamine: a hint that did not confirm}
Dopamine is cleared partly by a transporter. Removing that cleanup in the simulation changes how dopamine remains available, and we ask what, if anything, then changes in a specified downstream neural response. We use a history test: the network receives an earlier pattern and then a common final paired pattern, and the test asks whether the response to that final pattern depends on which pattern came before it. The comparison is between the normal transporter and a model of \textit{fumin}, whose transporter cleanup is disabled. In this test the patterned input is injected at TuBu, downstream of the eye; the fly does not see.

A second comparison asked whether lowering dopamine release altered the gap between the two transporter conditions. That manipulation is a fixed model multiplier, not a measured drug concentration or a treatment of a mutant fly. We first calibrated how the circuit distinguished the input patterns, then expressed the history contrast on that independent scale. We chose the primary comparisons, and their correction for testing more than one question, before examining outcomes. When the initial result looked promising but failed that correction, a separate follow-up fixed the genotype comparison and its decision rule in advance.

The first genotype contrast (\textit{fumin} minus normal, with unreduced release) was~\(\fixedthree{\result{GEN_MEAN}}\) template units, with a 95\% interval \roundedci{\result{GEN_CI}} and Holm-adjusted \(p=\fixedthree{\result{GEN_HOLM}}\). Its interval excluded zero, but the prespecified two-test correction did not support a finding. The change in that genotype gap under reduced release was +\(\fixedthree{\result{INTERACTION_MEAN}}\), interval \roundedmixedci{\result{INTERACTION_CI}}: no detected difference, and no established rescue.

In the independently frozen follow-up the genotype contrast was~\(\fixedthree{\result{CONF_MEAN}}\) template units (95\% interval \roundedmixedci{\result{CONF_CI}}; paired sign-randomisation \(p=\fpeval{round(\result{CONF_P},2)}\)). This is \emph{no detected difference}, not equivalence: the earlier negative hint was \result{CONF_DECISION} under the frozen criterion.

The follow-up is the more direct check of the initial hint, and it did not confirm it under the declared criterion. A confidence interval spanning zero leaves both directions of difference compatible with these runs, so removing the transporter could still matter. The comparison concerns this input, readout, model operating state and assumed receptor distribution. The unstable affine secondary calculation cannot identify an inhibitory mechanism, and the release comparison does not establish rescue. The runs repeat one simulated network with different random input, so their spread estimates the model's own noise. We do not infer an ADHD, attention or behavioural result.

\section*{Chemical knockout tour: a stress test of the model}
Before using the simulated brain for drug-like manipulations, we checked whether its activity reacts intelligibly when broad chemical pathways are removed. The model assigns many connections a transmitter label and an effective sign. We switched off the outgoing connections of each class in turn and left the network and its background drive otherwise in place. This is a deliberately blunt stress test: it removes transmission everywhere in a class at once, whereas a real antagonist's access and receptor targets would vary from cell to cell.

We compared each disconnected network with matched resting runs and concentrated on firing beyond the directly driven sensory neurons, because a sensory cell may still fire from its imposed background while its signal fails to propagate. The resting input, including photoreceptor background drive, is injected; the fly does not see. Figure~\ref{fig:tour} draws the model's changes on the anatomy.

Removing acetylcholine connections decreased firing outside sensory cells by 1.436 Hz [−1.451, −1.426], or 99.9\% of the 1.438 Hz control rate; removing GABA increased it by 12.588 Hz [+12.563, +12.613], and removing glutamate increased it by 10.779 Hz [+10.760, +10.799]. These directions fit the planned excitation/inhibition checks, but the GABA synchrony statistic did not move uniformly as predicted.

Histamine removal yielded no detected difference in optic-lobe mean firing (−0.001 Hz [−0.002, +0.001]), despite injected dark-rest photoreceptor drive. Dopamine disconnection likewise gave no detected outside-sensory mean difference (+0.001 Hz [−0.019, +0.025]). Removing octopamine connections lowered outside-sensory firing by 0.041 Hz [−0.060, −0.023], while serotonin removal raised it by 0.013 Hz [+0.002, +0.029]. Disconnecting neurons with no clear transmitter label reduced outside-sensory firing by 0.130 Hz [−0.159, −0.100].

The large changes after removing major excitatory or inhibitory classes are useful checks: altering the network does change its simulated activity. They do not validate the size of any response in a living fly, because a whole-class disconnection is much stronger and less selective than most pharmacological interventions. We keep the unresolved histamine and dopamine averages in view, since they could reflect this drive and readout rather than an absence of physiological function. The unlabelled-class result shows that gaps in transmitter annotation can have consequences. The tour can show where the model is sensitive and how it fails; whether every named chemical acts as expected, and what would happen in a seeing animal, remain open questions.

\section*{Courtship circuit: a signal reaches some downstream cells}
A brain-wide firing average cannot tell us whether activity travels along a particular proposed route, so we looked at a named courtship circuit. P1 cells are an entry point, pIP10 cells descend toward motor circuitry, and dPR1 and wing motor cells give more specific downstream readouts. If direct stimulation of P1 changes those readouts, the model has a route by which that imposed activity can influence them. It would not follow that a fly courts: cells can spike without a complete motor programme or a responsive body.

We applied different strengths of direct stimulation to P1 and also stimulated pIP10 directly. A matched comparison that stimulates random neurons instead asks whether a downstream change is specific to the selected entry point or would follow from adding drive anywhere in the graph. In these runs the signals are injected into the named neurons alongside background drive; the fly does not see. We tracked group firing and checked whether the proposed ordered chain of first spikes was present (Fig.~\ref{fig:courtship}).

Under strong direct P1 stimulation, P1 firing increased by 32.457 Hz [31.928, 32.837], pIP10 by 35.940 Hz [32.370, 38.420], dPR1 song-circuit cells by 26.230 Hz [23.390, 29.650] and wing motor cells by 9.159 Hz [8.203, 10.150], relative to matched controls. At medium input, dPR1's P1 response exceeded its matched random-neuron stimulation response by 17.810 Hz [15.810, 19.820]; direct pIP10 stimulation also raised dPR1 firing by 14.410 Hz [11.360, 17.830]. Yet vPR6 did not spike, and resting wing motor spikes make a simple first-spike causal sequence untenable.

A re-analysis found no demonstrated low-drive fighting versus high-drive song route: pIP10 responded even at low drive while the proposed pIP1 aggression route was not detected as a paired-cell response.

Activity reaching song-related and motor cells shows transmission through parts of the implemented circuit, and the random-neuron comparison makes the dPR1 result more informative than a stimulated-versus-rest comparison alone. But reaching a labelled cell does not prove the proposed sequence, and the missing vPR6 response matters alongside the cells that did respond. In particular, assigning low input to fighting and high input to song would oversell a split the re-analysis did not demonstrate. All of these changes are in circuit firing, and the body playback in the next section displays some of the same spikes.

\section*{GABA conductance-response: a steep network response}
A single all-or-nothing knockout cannot show how a circuit responds as an inhibitory brake is gradually weakened. Here we turn down, in steps, the strength of connections assigned to GABA, the model's major inhibitory class, and ask whether firing changes gently throughout or remains restrained until inhibition is almost gone. We also turn the GABA brake up, and separately strengthen GluCl, another inhibitory route, while GABA is partly blocked. If the second route brings activity toward rest, that is compensation within this network while the original GABA block stays in place.

Rdl homomers are sensitive to picrotoxin \cite{zhang}, but this model does not calibrate a drug concentration. The block axis shows the fraction of GABA conductance removed, not how much drug reaches cells in a male fly. We read out firing outside directly driven sensory cells so that a change in propagated activity is not confused with a change in the imposed source. The runs are at dark rest with background drive injected into sensory cells only and no TuBu stimulus; the fly does not see.

A \emph{half-effect point} is the input level at which a response reaches half a stated reference change; its meaning depends on whether the reference is an observed endpoint or an extrapolated maximum. \emph{Receptor reserve} means that a system can retain much of its response even after some receptors or their effective signalling have been lost. Here ``apparent reserve'' describes the network's firing curve, not a counted stock of spare Rdl receptors. A \emph{Hill slope} describes how sharply a fitted curve rises around its midpoint; it is a shape parameter, not by itself evidence for cooperative binding. We also check synchrony. The \emph{largest-1-ms-fraction} is the biggest share of the recorded population spiking together in any one-millisecond bin. The \emph{Fano factor} divides the variance of population spike counts across one-millisecond bins by their mean count, so a larger value means more variable counts relative to the average. Neither measure alone diagnoses a seizure.

The model's control outside-sensory rate was 1.438 Hz. At 10\%, 25\%, 50\%, 75\%, 90\% and full GABA block, the changes were respectively +0.093, +0.295, +0.972, +4.591, +9.591 and +12.588 Hz (Fig.~\ref{fig:dose}). Their paired 95\% intervals were [0.075, 0.107], [0.276, 0.318], [0.917, 1.022], [4.536, 4.652], [9.407, 9.753] and [12.563, 12.613] Hz. The rise is slow at first and steep near full block. Half the \emph{observed full-block change} occurs at 80.1\% blocked [79.9\%, 80.3\%]: the model has apparent functional reserve at low block. A freely fitted network slope was 5.05 [4.76, 5.34], but its inferred half-effect sat at the edge of the block range (99.9\% [96.1\%, 105.4\%]), with an interval running past full block. This is not an identified network EC50.

Of the frozen expectations, monotonic block and boost responses, a fitted slope above 1 and partial GluCl compensation matched; the predicted half-effect below 50\% block failed, the proposed central-brain-first order was opposite, and a sharp synchrony tip remained unclear. The spike-synchrony indicators did not reveal one clean seizure switch: the largest simultaneous-spike fraction rose, whereas the Fano statistic did not rise monotonically. Central-brain activity did not lead the optic lobes on their within-region half-response measure. Adding GABA conductance in steps of 0.25, 0.50, 1.00 and 2.00 times baseline lowered outside-sensory firing by 0.260 [0.238, 0.281], 0.442 [0.432, 0.456], 0.611 [0.598, 0.625] and 0.835 [0.823, 0.847] Hz, respectively; these are shallower changes than full block. Under 75\% GABA block, GluCl conductance increases of 0.50, 1.00 and 2.00 times baseline removed 23.5\% [20.9\%, 26.0\%], 45.5\% [44.7\%, 46.5\%] and 75.0\% [73.7\%, 76.1\%] of the excess firing. Even at the strongest reinforcement, the rate remained +1.147 Hz [1.103, 1.205] above rest.

Given the fitted curve's unstable half-effect, we report its slope only as a conditional fit, not a measured potency. Re-running the ladder on the GPU gave a fitted block Hill slope of 4.99 [4.65, 5.43] versus Brian2's 5.05 [4.76, 5.34], and half the observed full-block change at 80.12\% [79.95\%, 80.30\%] versus 80.11\% [79.90\%, 80.33\%] blocked; each pair of intervals overlaps. At the three GluCl reinforcements, GPU versus Brian2 removed 23.6\% [21.8\%, 25.2\%] versus 23.5\% [20.9\%, 26.0\%], 45.1\% [43.8\%, 46.2\%] versus 45.5\% [44.7\%, 46.5\%], and 75.0\% [74.0\%, 75.9\%] versus 75.0\% [73.7\%, 76.1\%] of excess firing; all three interval pairs overlap. The GPU boost rates had overlapping intervals with Brian2, but its fitted half-effect was not identified within the tested range (GPU $\phi=2.21$ [1.50, 4.29] versus Brian2 $\phi=1.26$ [0.96, 1.80]).

A second engine reproducing the main curve rules out an accident confined to one simulator's implementation as its sole cause. The shallow early block response and steeper late response tell us how this particular network absorbs a uniform loss of GABA conductance under the chosen drive. Because the fitted midpoint sits at the edge of the block range and its interval passes full block, the fitted network curve should not be read as a measured drug potency. Drug absorption, receptor composition and off-target effects could change the relation to a real fly entirely.

The two synchrony summaries answer different questions. The regional half-responses compare normalised changes within regions, not the arrival time of a seizure. The GluCl result is partial model rescue: raising a different inhibitory conductance can offset some extra firing while the GABA channel stays blocked. For pharmacology, a model dose curve should state what its axis means, show every sampled level and keep an observed half-change separate from an unstable fitted half-effect.

\section*{A virtual body replaying recorded spikes}
Spikes in wing motor cells are hard to picture as motion, so we made a video that plays recorded model spikes into NeuroMechFly v2 virtual fly bodies through FlyGym \cite{wangchen}, using the MuJoCo physics engine \cite{mujoco}. The video only illustrates a readout: the neural runs were completed before the bodies moved on screen, and their output cannot return as new sensory input. Input was injected directly into P1 and background sensory cells; the fly does not see.

We compared a playback from direct P1 stimulation with a matched resting playback. A chosen rule maps the activity of selected wing motor cells to wing angles. Choosing that rule is like choosing how to draw a graph: a different mapping could make the same spikes look like different movement. The comparison helps a reader see that the stimulation condition can be assigned a visibly different output while also showing whether the resting condition stays still.

The stimulated fly's wings (right of the clip) averaged about 47$^\circ$ during direct stimulation and about 14$^\circ$ afterward, versus about 8$^\circ$ for the resting body. The resting body's wings move too.

The angles come from a freely chosen spike-to-wing mapping, so the separation in the clip is a property of the recorded spikes and of our display rule. It does not reveal a calibrated motor transfer function or a working sensory-to-action loop. Watching the clip is useful for asking what a future body model would need to explain; it is not evidence of learned behaviour, a courtship song or real wing kinematics. The same caution applies to the static anatomical and activity figures, which make an implemented hypothesis easier to inspect.

\section*{Discussion: what the bench can tell us}
These studies show why a virtual pharmacology bench needs more than one striking image or one statistically suggestive result. The transporter experiment asked a narrow question defined in advance, and its follow-up did not detect the earlier hinted difference. In the knockout tour, broad interventions could strongly alter activity, and some transmitter readouts stayed unresolved. Courtship stimulation reached named downstream cells without the proposed behavioural split. The GABA series gave the most directly interpretable graded response, including a partial compensatory route, although its blocked-fraction axis is an assumed uniform model manipulation, not a drug dose. Playing spikes into the virtual body made an output visible without turning any of those neural observations into behaviour. Together they are complementary checks on one constructed system, not independent biological replications.

For a pharmacology student, the dopamine null is a result worth reporting. A result with an interval spanning zero says the declared comparison did not resolve a difference under these conditions. The original hint is unconfirmed, and dopamine clearance may still matter. The next question is whether the input, state and readout could detect an effect that matters biologically, using criteria chosen before seeing the answer. Searching other readouts after the fact could at most generate hypotheses. Because the receptor map is assumed, even a clean positive genotype contrast would first describe the behaviour of the model's dopamine implementation.

The knockout tour works like a test of the controls on laboratory equipment. When removing an excitatory class suppresses network firing, or removing an inhibitory class increases it, we learn that the implemented wiring and signs have consequences. A surprising null is a prompt to inspect how that class is driven and which population was measured. None of these disconnections has the selectivity, access or graded exposure of a drug. The courtship result likewise shows why an anatomically plausible path and increased firing at its endpoints are not enough to claim a motor programme.

The dose curve asks how much the measured network output changes as an explicitly defined model target is varied. Its early flatness and late steepness are worth explaining, but the half-effect must name its reference, and a fitted Hill slope cannot turn an extrapolated midpoint into a measured potency. The GluCl compensation is a testable prediction of interactions between implemented brakes. Before relating it to treatment, one would need native receptor sensitivity, distribution and independent physiological readouts. A drug that binds one target in a dish can reach several targets in an animal, with an entirely different whole-organism dose response.

Several unvalidated translations separate a claim about this model from a claim about flies. MaleCNS supplies fly anatomy \cite{berg}; the assigned electrical parameters, transmitter effects and largely uniform dopamine receptor densities supply a computational operating point. Shiu and colleagues' experimental precedent \cite{shiu} does not validate this particular operating point. Injected input bypasses natural sensory transformation, and the body is one-way playback. Matching the GPU engine to Brian2 checks that two implementations compute the tested spike outputs consistently, and those biological gaps remain. The next step is an independently constrained prediction that could fail when compared with a living fly.

\section*{Limitations and next steps}
One anatomical specimen and assumed electrical parameters are repeated with random inputs: uncertainty intervals do not measure between-fly variability or uncertainty in the model's receptor map. The EM connectome does not resolve electrical synapses (gap junctions), so the model has none. Neuropeptide co-transmission is absent, and dopamine is the only neuromodulator with modelled dynamics; other labelled transmitters, including serotonin and octopamine, contribute only assigned synaptic signs. All point neurons share the same base electrical parameters. At rest, background input is injected into sensory cells only. An inhibitory synapse onto a cell that never fires cannot reduce its firing further, so blocking GABA cannot reveal disinhibition in parts of the network that the injected input does not reach. Uniform transmitter switches and receptor multipliers omit native cell-specific pharmacology. The optic lobe and visual projection neurons are nearly silent at rest in this model (0.023 and 0.059 Hz versus 4.795 Hz centrally); the outside-sensory mean is therefore driven mainly by the central brain, and the histamine result depends on this resting state. The virtual body is one-way playback; the GPU equivalence result is limited to its tested inputs and clamped dopamine. No result establishes behaviour or clinical relevance. A next experiment needs labelled receptor and cell-type uncertainty, measured circuit constraints, validated sensory input and an agreed prospective prediction before using the model for a biological claim.

\section*{Data and code availability}
The anatomical and synaptic connectivity data derive from the MaleCNS v1.0 connectome, released under CC BY 4.0 by the FlyEM Project Team at HHMI Janelia, the University of Cambridge, the MRC Laboratory of Molecular Biology, and Google Research. The dataset is available at \url{https://male-cns.janelia.org/}, \url{gs://flyem-male-cns} and \url{https://neuprint.janelia.org/}. The model code, reviewed simulation summaries and figure sources are publicly available under the MIT License at \url{https://github.com/uguryildirim24/flyonenomics} (release v1.0.0). The v1.0.0 release is archived on Zenodo (DOI: \url{https://doi.org/10.5281/zenodo.23070186}).

\section*{Author contributions and AI use}
Hasan ``Rolf'' Yildirim conceived the question, directed the project and is responsible for the content. The code, analyses and text were produced with AI assistance: Claude (Anthropic) drafted and reviewed text and code, OpenAI Codex implemented and checked code and analyses, and Gemini (Google) checked literature and citations. These tools are not authors, and their checks are not independent human or experimental replication. \iffinal\else Rolf must approve this draft and disclosure before submission.\fi

\section*{Competing interests}
The author declares no competing interests.

\section*{Funding}
This work received no external funding.

\begin{figure}[p]\centering
\includegraphics[width=\textwidth]{../../figures/3d/male-cns-anatomy-paper.png}
\caption{Static anatomical context from MaleCNS. Coloured structures locate representative groups. MaleCNS v1.0, Berg et al. \cite{berg}, CC BY 4.0.}\label{fig:anatomy}
\end{figure}
\begin{figure}[p]\centering
\includegraphics[width=\textwidth]{../../figures/tour/knockout-overview.png}
\caption{Anatomical rate-change illustration of transmitter disconnections, with the same colour range across panels. The shared colour scale marks $\pm200$ and $\pm10$ Hz; GABA and glutamate changes can saturate the display. Displayed skeletons sample the graph; the analysis uses all retained cells. The input is injected, and the disconnections are made in the model.}\label{fig:tour}
\end{figure}
\begin{figure}[p]\centering
\includegraphics[width=\textwidth]{../../figures/tour/courtship-path.png}
\caption{Courtship-path anatomy and model rate changes under imposed P1 drive. Highlighted routes are not evidence of movement or song.}\label{fig:courtship}
\end{figure}
\begin{figure}[p]\centering
\includegraphics[width=\textwidth]{gaba-curves.pdf}
\caption{GABA response curves: outside-sensory firing differences from matched model rest, with whole-repeat intervals. Horizontal axes show the fraction of GABA conductance blocked or conductance multipliers, not administered dose.}\label{fig:dose}
\end{figure}
\begin{thebibliography}{18}\small\raggedright
\bibitem{holm} Holm S (1979). A simple sequentially rejective multiple test procedure. \textit{Scandinavian Journal of Statistics} 6:65--70.
\bibitem{efron} Efron B, Tibshirani RJ (1993). \textit{An Introduction to the Bootstrap}. Chapman \& Hall/CRC.
\bibitem{ernst} Ernst MD (2004). Permutation methods: a basis for exact inference. \textit{Statistical Science} 19(4).
\bibitem{kume} Kume K, et al. (2005). Dopamine is a regulator of arousal in the fruit fly. \textit{Journal of Neuroscience}. \url{https://doi.org/10.1523/JNEUROSCI.2048-05.2005}.
\bibitem{shin} Shin M, Venton BJ (2018). Electrochemical Measurements of Acetylcholine-Stimulated Dopamine Release in Adult Drosophila melanogaster Brains. \textit{Analytical Chemistry}. \url{https://doi.org/10.1021/acs.analchem.8b02114}.
\bibitem{berg} Berg S, et al. (2026). Sexual dimorphism in the complete Drosophila male central nervous system connectome. \textit{Cell}. \url{https://doi.org/10.1016/j.cell.2026.08.015}.
\bibitem{shiu} Shiu PK, et al. (2024). A Drosophila computational brain model reveals sensorimotor processing. \textit{Nature}. \url{https://doi.org/10.1038/s41586-024-07763-9}.
\bibitem{zhang} Zhang HG, Lee HJ, Rocheleau T, ffrench-Constant RH, Jackson MB (1995). Subunit composition determines picrotoxin and bicuculline sensitivity of Drosophila gamma-aminobutyric acid receptors. \textit{Molecular Pharmacology} 48(5):835--840. PMID 7476913.
\bibitem{wangchen} Wang-Chen S, et al. (2024). NeuroMechFly v2: simulating embodied sensorimotor control in adult Drosophila. \textit{Nature Methods} 21(12):2353--2362. \url{https://doi.org/10.1038/s41592-024-02497-y}.
\bibitem{mujoco} Todorov E, Erez T, Tassa Y (2012). MuJoCo: A physics engine for model-based control. \textit{2012 IEEE/RSJ International Conference on Intelligent Robots and Systems}, 5026--5033. \url{https://doi.org/10.1109/IROS.2012.6386109}.
\bibitem{stimberg} Stimberg M, Brette R, Goodman DFM (2019). Brian 2, an intuitive and efficient neural simulator. \textit{eLife} 8:e47314. \url{https://doi.org/10.7554/eLife.47314}.
\bibitem{salmon} Salmon JK, Moraes MA, Dror RO, Shaw DE (2011). Parallel random numbers: as easy as 1, 2, 3. \textit{SC '11}, article 16:1--12. \url{https://doi.org/10.1145/2063384.2063405}.
\bibitem{okuta} Okuta R, Unno Y, Nishino D, Hido S, Loomis C (2017). CuPy: A NumPy-compatible library for NVIDIA GPU calculations. \textit{NIPS LearningSys Workshop}, paper 16. \url{http://learningsys.org/nips17/assets/papers/paper_16.pdf}.
\bibitem{nvrtc} NVIDIA Corporation (2024). NVIDIA CUDA Runtime Compilation (NVRTC) Library. \textit{CUDA Toolkit Documentation}, version 12.6. \url{https://docs.nvidia.com/cuda/nvrtc/index.html}.
\bibitem{alevi} Alevi D, et al. (2022). Brian2CUDA: Flexible and efficient simulation of spiking neural network models on GPUs. \textit{Frontiers in Neuroinformatics} 16:883700. \url{https://doi.org/10.3389/fninf.2022.883700}.
\bibitem{knight} Knight JC, Nowotny T (2018). GPUs outperform current HPC and neuromorphic solutions in terms of speed and energy when simulating a highly-connected cortical model. \textit{Frontiers in Neuroscience} 12:941. \url{https://doi.org/10.3389/fnins.2018.00941}.
\end{thebibliography}
\end{document}
